% 数学ノート集・共通テンプレート v2（積分論2026年度第1部の版面を参考）
% A5判。章内で定義・定理・例・注意等の番号を共通化。
% コンパイル： platex -interaction=nonstopmode math_note_template_v2.tex (2回)
%              dvipdfmx math_note_template_v2.dvi
\documentclass[a5paper,10pt,onecolumn]{jsarticle}
\usepackage[dvipdfmx,top=22mm,bottom=20mm,inner=19mm,outer=18mm,headheight=13pt,headsep=6mm]{geometry}
\usepackage{amsmath,amssymb,amsthm,mathtools}
\usepackage{enumitem,fancyhdr,titlesec}
\usepackage[dvipdfmx,bookmarks=true,colorlinks=false]{hyperref}
\usepackage{pxjahyper}

% ----- 差し替える文書情報 -----
\newcommand{\NoteTitle}{Riemann計量とLevi-Civita接続}
\newcommand{\NoteSeries}{数学ノート集}

% ----- 文字・組版 -----
% 参考PDFの質素な明朝＋Computer Modernの印象を引き継ぎながら，
% A5判で窮屈にならないよう本文の行間を少し確保する．
\linespread{1.13}
\setlength{\parindent}{1zw}
\setlength{\parskip}{0pt}
\setlength{\abovedisplayskip}{6pt plus 2pt minus 2pt}
\setlength{\belowdisplayskip}{6pt plus 2pt minus 2pt}
\setlength{\abovedisplayshortskip}{3pt plus 2pt minus 1pt}
\setlength{\belowdisplayshortskip}{4pt plus 2pt minus 1pt}
\setlist[enumerate]{topsep=3pt,itemsep=1pt,parsep=0pt,leftmargin=2.4zw}
\titleformat{\section}{\normalfont\normalsize\bfseries}{\thesection}{0.8zw}{}
\titlespacing*{\section}{0pt}{2.0ex plus .4ex}{.9ex}
\titleformat{\subsection}{\normalfont\normalsize\bfseries}{\thesubsection}{0.7zw}{}
\titlespacing*{\subsection}{0pt}{1.5ex plus .3ex}{.6ex}
\raggedbottom

% ----- 柱とページ番号：ページ上部、罫線なし -----
\fancypagestyle{mathnote}{%
  \fancyhf{}%
  \fancyhead[L]{\footnotesize \NoteSeries\ （\NoteTitle）}%
  \fancyhead[R]{\footnotesize \thepage}%
  \renewcommand{\headrulewidth}{0pt}%
  \renewcommand{\footrulewidth}{0pt}%
}
\pagestyle{mathnote}

% ----- 節内通し番号の定義・定理群 -----
\newtheoremstyle{mathnotetheorem}
 {4.5pt}{4.5pt}{\normalfont}{}%
 {\bfseries}{}{0.65zw}%
 {\thmname{#1}\thmnumber{ #2}\thmnote{（#3）}}
\theoremstyle{mathnotetheorem}
\newtheorem{theorem}{定理}[section]
\newtheorem{definition}[theorem]{定義}
\newtheorem{lemma}[theorem]{補題}
\newtheorem{proposition}[theorem]{命題}
\newtheorem{corollary}[theorem]{系}
\newtheorem{example}[theorem]{例}
\newtheorem{remark}[theorem]{注意}
\renewcommand{\proofname}{証明}
\renewcommand{\qedsymbol}{$\square$}
\makeatletter
\renewenvironment{proof}[1][\proofname]{%
  \par\pushQED{\qed}\normalfont\topsep4\p@\@plus2\p@\relax
  \trivlist\item[\hskip\labelsep\bfseries #1]\ignorespaces
}{\popQED\endtrivlist\@endpefalse}
\makeatother

% ----- 使うノートで適宜拡張する記号 -----
\newcommand{\R}{\mathbb{R}}
\newcommand{\norm}[1]{\lVert #1\rVert}
\newcommand{\abs}[1]{\lvert #1\rvert}

\begin{document}
\section*{読了後にできること}
本稿を読むための前提は，多様体の局所座標とChristoffel記号，Koszul公式．
本稿の具体的な到達課題は，計量の係数からLevi--Civita接続のChristoffel記号を算出する．


\section{Riemann計量}
\begin{definition}
滑らかな多様体 $M$ のRiemann計量とは，
各 $p\in M$ で正定値内積を与える
滑らかな対称 $(0,2)$ テンソル場 $g$ をいう．
$g(X,Y)$ を $\langle X,Y\rangle$ とも書く．
\end{definition}
\begin{theorem}[計量の存在]\label{thm:exists}
Hausdorff第二可算滑らかな多様体はRiemann計量を持つ．
\end{theorem}
\begin{proof}
座標近傍 $(U_i,x_i)$ 上で
$g_i=\sum_jdx_i^j\otimes dx_i^j$ と定め，
被覆に従属する局所有限な滑らかな単位の分割 $(\rho_i)$ を取る．
\[
 g_p=\sum_i\rho_i(p)(g_i)_p
\]
と置く．局所有限性から滑らかで対称である．
零でない $v\in T_pM$ に対し
少なくとも一つの $\rho_i(p)>0$ があり，
各 $g_i(v,v)>0$ なので $g_p(v,v)>0$．
従って計量となる．
\end{proof}
\section{接続の捩率と計量保存}
\begin{definition}
$TM$ 上の接続 $\nabla$ の捩率を
\[
 T(X,Y)=\nabla_XY-\nabla_YX-[X,Y]
\]
と定める．
捩率零の接続を対称接続という．
$\nabla$ が計量と両立するとは
\[
 Xg(Y,Z)=g(\nabla_XY,Z)+g(Y,\nabla_XZ)
\]
が全てのベクトル場について成立することをいう．
\end{definition}
\begin{lemma}[捩率のテンソル性]
$T$ は $(1,2)$ 型テンソル場である．
\end{lemma}
\begin{proof}
$T(fX,Y)$ を展開し，
$\nabla_Y(fX)=Y(f)X+f\nabla_YX$ と
$[fX,Y]=f[X,Y]-Y(f)X$ を使えば
$T(fX,Y)=fT(X,Y)$．
反対称性から第2変数も同様である．
\end{proof}
\section{Koszul公式}
\begin{theorem}[Levi-Civitaの存在と一意性]\label{thm:levi}
任意のRiemann計量 $g$ に対して，
捩率零かつ計量と両立する接続 $\nabla$ が
ただ一つ存在する．これは
\[
\begin{aligned}
2g(\nabla_XY,Z)
={}&Xg(Y,Z)+Yg(Z,X)-Zg(X,Y)\\
&-g(X,[Y,Z])+g(Y,[Z,X])+g(Z,[X,Y])
\end{aligned}
\tag{K}
\]
によって定まる．
\end{theorem}
\begin{proof}
まず存在するなら式 (K) が必要であることを示す．
計量保存の式を $(X,Y,Z)$，
$(Y,Z,X)$，$(Z,X,Y)$ の順に書き，
最初の二式から三番目を引く．
$\nabla_XY-\nabla_YX=[X,Y]$ を代入して
残る項を整理すると (K) を得る．
従って接続は一意である．

逆に (K) の右辺を $2B(X,Y,Z)$ とする．
$Z$ について滑らかな関数 $f$ を掛けた場合，
Lie括弧のLeibniz則
$[Y,fZ]=f[Y,Z]+Y(f)Z$ 等を使うと，
$f$ の導関数を含む項が相殺し，
$B(X,Y,fZ)=fB(X,Y,Z)$ を得る．
よって内積の非退化性から，任意の $X,Y$ に対し
$g(\nabla_XY,Z)=B(X,Y,Z)$ を満たす
ただ一つの滑らかなベクトル場 $\nabla_XY$ が存在する．
同様の展開によって
$B(fX,Y,Z)=fB(X,Y,Z)$ と
\[
 B(X,fY,Z)=X(f)g(Y,Z)+fB(X,Y,Z)
\]
を得る．
したがって $\nabla$ は接続である．
式 (K) において $X,Y$ を交換して差を取ると
$2g(\nabla_XY-\nabla_YX,Z)=2g([X,Y],Z)$，
よって捩率零である．
また $Y,Z$ を交換して和を取ると
$g(\nabla_XY,Z)+g(Y,\nabla_XZ)=Xg(Y,Z)$ となるので
計量を保存する．
\end{proof}
\section{局所座標による表示}
\begin{proposition}[Christoffel記号]\label{prop:christoffel}
座標 $x^i$ において
$\nabla_{\partial_i}\partial_j
=\sum_k\Gamma^k_{ij}\partial_k$ と書き
$g_{ij}=g(\partial_i,\partial_j)$ とすると，
\[
 \Gamma^k_{ij}=
 \frac12\sum_\ell g^{k\ell}
 (\partial_i g_{j\ell}+\partial_jg_{i\ell}-\partial_\ell g_{ij})
\]
が成立する．
\end{proposition}
\begin{proof}
座標ベクトル場のLie括弧は零なのでKoszul公式から
$2\sum_kg_{k\ell}\Gamma^k_{ij}
=\partial_i g_{j\ell}+\partial_jg_{i\ell}
-\partial_\ell g_{ij}$．
逆行列 $(g^{k\ell})=(g_{ij})^{-1}$ を掛けて解けばよい．
\end{proof}
\begin{proposition}[計量に沿った長さの保存]
$\gamma$ に沿って $\nabla_{\dot\gamma}V=0$，
$\nabla_{\dot\gamma}W=0$ を満たすベクトル場 $V,W$ に対し
$g(V,W)$ は $\gamma$ 上一定である．
\end{proposition}
\begin{proof}
計量保存性を曲線に沿って適用すると
$\frac d{dt}g(V,W)=
g(\nabla_{\dot\gamma}V,W)+g(V,\nabla_{\dot\gamma}W)=0$ である．
\end{proof}

\section{到達課題の完全解答}
\noindent\textbf{極座標の計量から接続を復元する．}
\begin{proposition}[Euclid平面の極座標]
$\mathbb R^2\setminus\{0\}$ の極座標 $(r,\theta)$
におけるEuclid計量は
$g=dr^2+r^2d\theta^2$ である．
非零のChristoffel記号は
\[
\Gamma^r_{\theta\theta}=-r,\qquad
\Gamma^\theta_{r\theta}
=\Gamma^\theta_{\theta r}=\frac1r
\]
のみである．
\end{proposition}
\begin{proof}
一般公式
\[
\Gamma^k_{ij}
=\frac12\sum_\ell g^{k\ell}
(\partial_i g_{j\ell}+\partial_j g_{i\ell}
-\partial_\ell g_{ij})
\]
を用いる．この場合
$(g_{ij})=\operatorname{diag}(1,r^2)$，
$(g^{ij})=\operatorname{diag}(1,r^{-2})$ であり，
非零の計量係数の導関数は
$\partial_r g_{\theta\theta}=2r$ のみである．
代入すると表示式となる．
\end{proof}
\begin{corollary}[極座標での測地線方程式]
測地線は
\[
\ddot r-r\dot\theta^2=0,\qquad
\ddot\theta+\frac2r\dot r\dot\theta=0
\]
を満たす．第二式は
$(r^2\dot\theta)'=0$ と同値である．
\end{corollary}
\begin{proof}
測地線方程式に上のChristoffel記号を代入する．
第二式に $r^2$ を掛けて積の微分を取れば
$(r^2\dot\theta)'=0$ が得られる．
\end{proof}
\noindent\textbf{到達確認．}
$\theta$ が一定のとき，測地線方程式を解け．
\par\noindent\textbf{解答．}
$\dot\theta=0$ なので $\ddot r=0$ であり，
$r(t)=at+b>0$ となる定義区間で径方向の直線を表す．
\end{document}
